% Hướng dẫn sử dụng phòng Kinh doanh ERP Sao Việt Nhật: việc hằng ngày của Trợ lý kinh doanh và Sale.
% Biên dịch: xelatex huong-dan-su-dung-kinh-doanh.tex (chạy hai lần cho mục lục).
% Phần khai báo (quyền, danh mục, khách, chính sách tài chính) nằm ở ../huong-dan-cau-hinh/cau-hinh-he-thong.tex.
\documentclass[11pt,a4paper]{article}
\usepackage{fontspec}
\usepackage[vietnamese]{babel}
\setmainfont{Times New Roman}
\newfontfamily\tieude{Arial}
\usepackage[a4paper,margin=1.9cm,top=1.7cm]{geometry}
\usepackage{graphicx}
\usepackage{enumitem}
\usepackage{xcolor}
\usepackage{titlesec}
\usepackage{tocloft}
\usepackage{needspace}
\usepackage{fancyhdr}
\usepackage[hidelinks]{hyperref}
\usepackage[most]{tcolorbox}
\usepackage{booktabs}
\usepackage{tabularx}
\graphicspath{{hinh/}}

\definecolor{gach}{HTML}{C2410C}
\definecolor{xamnen}{HTML}{F4F4F5}
\definecolor{xamchu}{HTML}{52525B}

\setlist{nosep,leftmargin=1.5em}
\setlength{\parskip}{0.4em}
\setlength{\parindent}{0pt}
\setlength{\emergencystretch}{2em}

\titleformat{\section}{\tieude\Large}{Chương \thesection.}{0.4em}{}[{\color{black!25}\titlerule[0.4pt]}]
\titlespacing*{\section}{0pt}{1.2em}{0.6em}
\titleformat{\subsection}{\tieude\large}{\thesubsection}{0.6em}{}
\titlespacing*{\subsection}{0pt}{1.1em}{0.3em}
\renewcommand{\cftsecpresnum}{Chương }
\setlength{\cftsecnumwidth}{5.2em}
\renewcommand{\cftsecfont}{\tieude}
\renewcommand{\cftsecpagefont}{\tieude}
\setlength{\cftbeforesecskip}{0.5em}
\addto\captionsvietnamese{\renewcommand{\contentsname}{Mục lục}}
\renewcommand{\cfttoctitlefont}{\tieude\large}
\setlength{\cftbeforetoctitleskip}{0pt}
\setlength{\cftaftertoctitleskip}{0pt}

\pagestyle{fancy}
\fancyhf{}
\fancyhead[L]{\tieude\footnotesize\color{xamchu}Hướng dẫn sử dụng phòng Kinh doanh}
\fancyhead[R]{\tieude\footnotesize\color{xamchu}ERP Sao Việt Nhật}
\fancyfoot[C]{\tieude\footnotesize\thepage}
\renewcommand{\headrulewidth}{0.4pt}

\newtcolorbox{boicanh}{colback=xamnen,colframe=xamnen,boxrule=0pt,arc=3pt,
  left=9pt,right=9pt,top=6pt,bottom=6pt,before skip=4pt,after skip=8pt}

\definecolor{vang}{HTML}{FEF08A}
\newcommand{\oo}[1]{{\setlength{\fboxsep}{1.2pt}\colorbox{vang}{\strut #1}}}
\newcommand{\duong}[1]{{\tieude\small #1}}
\newcommand{\so}[1]{{\tieude\bfseries\small\textcolor{white}{\colorbox{red!80!black}{\kern-2pt#1\kern-2pt}}}}
\newcommand{\khung}[2]{{\setlength{\fboxsep}{0pt}\setlength{\fboxrule}{0.5pt}%
  \fcolorbox{black!20}{white}{\includegraphics[width=\dimexpr#1-1pt\relax]{#2}}}}
\newcommand{\q}[1]{\textquotedbl#1\textquotedbl{}}
\newcommand{\canh}[3]{%
  \par\noindent
  \begin{minipage}[t]{#1\linewidth}\vspace{0pt}\khung{\linewidth}{#2}\end{minipage}\hfill
  \begin{minipage}[t]{\dimexpr0.97\linewidth-#1\linewidth\relax}\vspace{0pt}\setlength{\parskip}{0.5em}#3\end{minipage}\par\medskip}
\newcommand{\rong}[2][\linewidth]{\par\noindent\begin{center}\khung{#1}{#2}\end{center}\vspace{-0.6em}}
% Tên tài liệu khác.
\newcommand{\tl}[1]{\textit{#1}}
% Liên kết sang một mục của tài liệu cấu hình (đích là tên \label bên đó, nhờ destlabel).
\newcommand{\cauhinh}{../huong-dan-cau-hinh/cau-hinh-he-thong.pdf}
\definecolor{lienket}{HTML}{1E40AF}
\newcommand{\kbn}[2]{\href{\cauhinh\##1}{\textcolor{lienket}{#2}}}
\newcommand{\kb}[2]{\tl{Hướng dẫn cấu hình hệ thống}, mục \kbn{#1}{#2}}

\begin{document}

\thispagestyle{empty}
{\tieude\small\color{xamchu} ERP Sao Việt Nhật}\par\vspace{3pt}
{\tieude\fontsize{22}{26}\selectfont Hướng dẫn sử dụng phòng Kinh doanh}\par\vspace{4pt}
{\small Tính giá, báo giá, đơn hàng bán, theo dõi sản xuất, xuất hàng, tiền và hoa hồng}\par\vspace{2pt}
{\small\color{xamchu} Cập nhật 04/11/2026}\par\vspace{6pt}
\noindent{\color{black!25}\rule{\linewidth}{0.4pt}}\par\vspace{4pt}

Tài liệu cho Trợ lý kinh doanh và Sale, dùng hằng ngày. Đi theo đúng đường một đơn hàng: khách hỏi giá, lập phiếu
tính giá, làm báo giá, khách chốt, lên đơn, theo dõi sản xuất, xuất hàng, rồi tiền về và hoa hồng.

Phần khai báo trước khi dùng (quyền từng màn, danh mục giấy, máy, công đoạn, khách hàng, chính sách tài chính của
khách) nằm ở tài liệu \tl{Hướng dẫn cấu hình hệ thống}, chương \kbn{danh-muc}{Danh mục dùng chung} và chương
\kbn{kinh-doanh}{Kinh doanh}. Chữ màu xanh là liên kết, bấm vào mở đúng mục bên tài liệu đó.

Ảnh chụp từ phần mềm thật. Ô đỏ đánh số là chỗ cần bấm hoặc cần đọc, số khớp với lời giải thích bên cạnh.

\bigskip
\tableofcontents

% ============================================================
\Needspace*{16\baselineskip}
\section{Đường đi của một đơn hàng}\label{duong-di}

Sơ đồ dưới gồm hai phần. Hai hàng trên cùng là những thứ phải khai báo trước thì phiếu tính giá mới ra số. Các
hàng còn lại là đường đi của một đơn hàng, mỗi hàng là một vai hoặc một phòng. Mỗi ô là một bước, chữ trên mũi
tên là điều kiện hoặc thứ được chuyển sang bước sau.

\rong[0.86\linewidth]{kd-00-luong.png}

\Needspace*{12\baselineskip}
\subsection{Khai báo trước khi tính giá}\label{khai-bao-truoc}

Phần này do quản trị phần mềm hoặc người được giao khai danh mục làm, thường một lần lúc bắt đầu và mỗi khi có
giấy, máy, công đoạn mới. Trợ lý và sale không cần tự khai, nhưng cần biết số trên phiếu lấy từ đâu để báo đúng
người khi giá ra lạ. Thứ tự khai theo mũi tên, mỗi bước xem ở \tl{Hướng dẫn cấu hình hệ thống}:

\begin{enumerate}
\item \kbn{danh-muc-don-vi}{Đơn vị \& quy đổi}: tờ, kg, con, cái và cách đổi qua lại. Máy, giấy, vật tư đều cần
đơn vị.
\item \kbn{danh-muc-giay}{Giấy}: định lượng, giá một kg, công thức tính tiền giấy.
\item \kbn{danh-muc-may}{Thiết bị \& Máy móc}: khổ, tốc độ, phần chừa nhíp. Phiếu tính giá lấy khổ tờ in từ đây.
\item \kbn{danh-muc-vat-tu}{Vật tư khác}: mực, kẽm, keo, màng, mỗi thứ có công thức tính tiền và công thức định
mức.
\item \kbn{danh-muc-cong-doan}{Công đoạn}: nối máy và vật tư thành từng bước In, cán, bế, dán; có công thức tính
giá và bảng bù hao. Phiếu tính giá ghép các công đoạn thành chuỗi bắt đầu từ In.
\item \kbn{kinh-doanh-khach}{Khách hàng}: sale tự khai khách của mình, kèm
\kbn{kinh-doanh-tai-chinh}{chính sách tài chính} (hạn mức nợ, rào giá).
\end{enumerate}

Hai danh mục không cần khai trước. \kbn{danh-muc-khuon}{Khuôn bế}: phiếu tính giá chỉ hỏi có sẵn hay làm mới, khuôn
cụ thể do Kế hoạch sản xuất chọn khi lên lệnh. \kbn{danh-muc-thanh-pham}{Thành phẩm}: phần mềm tự tạo khi chốt đơn
hàng bán.

Một việc không cần cho tính giá nhưng phải xong trước khi đơn đầu tiên xuống xưởng: có người giữ quyền
\oo{Thao tác} ở màn KCS, xem \kb{danh-muc-kcs-quyen}{Ai được kiểm KCS}. Thành phẩm chỉ vào kho khi KCS kiểm xong và
gửi sang. Chưa ai có quyền này thì đơn đứng mãi ở dòng \q{Chờ KCS gửi thành phẩm sang kho} và không xuất hàng
được. \kbn{danh-muc-kcs}{Tiêu chí KCS} thì không bắt buộc, thiếu tiêu chí KCS vẫn kiểm được.

Khai xong thì lập thử một phiếu mẫu và so với giá tính tay, theo \kb{kinh-doanh-kiem}{Kiểm bằng một phiếu tính giá
mẫu}.

\Needspace*{12\baselineskip}
\subsection{Đường đi của đơn hàng}\label{duong-di-don}

Khách hỏi giá thì sale ghi yêu cầu vào hồ sơ khách. Trợ lý lập phiếu tính giá, rồi làm báo giá. Sale gửi báo giá
cho khách, ghi đồng ý hay từ chối, lên đơn và chốt đơn (trợ lý cũng lập đơn hộ được). Đơn cần cọc thì chờ Kế toán
thu đủ cọc mới sang sản xuất, đơn không cần cọc thì sang thẳng. Trong lúc sản xuất, trợ lý theo dõi tiến độ để
trả lời khách. Hàng vào kho thì sale bấm Xuất hàng ở đơn, thủ kho lập phiếu xuất và ghi sổ. Kho xuất xong, Kế toán
ghi hoá đơn và thu tiền, tiền về thì tính hoa hồng cho sale phụ trách khách.

Trên phần mềm, cùng luồng này nằm ở \duong{Kinh doanh > Quy trình kinh doanh}. Mỗi ô có số thứ tự, cột là phòng
làm bước đó, bấm vào một ô để mở đúng màn của bước.

\rong[0.8\linewidth]{kd-01-quy-trinh.png}

Phòng Kinh doanh làm sáu bước ở cột đầu: 1 ghi nhận khách, 2 lập phiếu tính giá, 3 làm báo giá, 4 gửi khách và ghi
kết quả, 5 lên đơn và chốt đơn, 13 xuất hàng. Các bước còn lại do phòng khác làm, phòng Kinh doanh chỉ theo
dõi để trả lời khách:

\begin{itemize}
\item Kế toán thu cọc (bước 6), ghi hoá đơn (15), thu tiền (16).
\item Kế hoạch sản xuất lên lệnh, lo vật tư, xếp lịch máy (7 đến 9). Tổ sản xuất chạy công đoạn, KCS kiểm hàng,
kho nhận thành phẩm (10 đến 12).
\item Kho soạn hàng, lập phiếu xuất và ghi sổ (14).
\end{itemize}

\Needspace*{10\baselineskip}
\subsection{Trợ lý và Sale khác nhau ở đâu}\label{vai}

Trợ lý kinh doanh lập phiếu tính giá, báo giá, đơn hộ cho mọi sale và thấy khách của tất cả sale. Trợ lý thấy từng
dòng giá vốn và công thức trên phiếu.

Sale làm với khách mình phụ trách: thêm khách, xem báo giá, gửi khách, ghi kết quả, lên đơn, theo dõi đơn. Sale
không thấy chi tiết giá vốn và công thức, trên phiếu tính giá chỉ thấy giá vốn tổng.

Báo giá và đơn luôn đứng tên người phụ trách khách, kể cả khi trợ lý lập hộ. Ví dụ ở chương này: Admin lập báo giá
cho khách CAFÉ KATINAT, nhưng báo giá ghi \oo{NV phụ trách Nguyễn Thị Huyên}, vì chị Huyên phụ trách khách đó.
Người phụ trách cũng là người hưởng hoa hồng khi tiền về (chương \ref{tien}).

Quyền cụ thể từng vai do quản trị phần mềm cài, xem \kb{kinh-doanh-quyen}{Quyền của sáu màn}. Bảng và nút nào bạn không thấy trong tài liệu này là do vai của bạn chưa được bật.

% ============================================================
\Needspace*{16\baselineskip}
\section{Khách hàng}\label{khach}

\begin{boicanh}
Khách mới thì tạo hồ sơ trước khi tính giá (cách tạo ở \kb{kinh-doanh-khach}{Khai khách hàng}). Khách cũ thì mở hồ sơ để xem đã mua gì, báo giá cũ ra sao, rồi lập phiếu tính giá ngay từ hồ sơ.
\end{boicanh}

\duong{Kinh doanh > Khách hàng}. Gõ tên, mã số thuế hoặc số điện thoại vào ô tìm. Bấm vào tên khách để mở hồ sơ.

\rong[0.8\linewidth]{kd-60-khach.png}

\so{1} \oo{Lập phiếu tính giá}: mở phiếu mới đã chọn sẵn khách này.

\so{2} \oo{Lịch sử mua hàng} và \oo{Lịch sử báo giá}: các đơn và báo giá cũ của khách, bấm một dòng để mở. Khách
hỏi lại giá hàng đã làm thì xem ở đây trước.

\so{3} Số liệu của khách trong kỳ chọn ở hàng trên: doanh số, số đơn đã chốt, nhịp đặt hàng, tỉ lệ chốt báo giá.

Các tab còn lại: \oo{Chăm sóc} để hẹn gọi, gặp khách, \oo{Liên hệ} là người liên hệ của
khách, \oo{Giao hàng} là các địa chỉ nhận hàng. Cuối tab \oo{Tổng quan} có thẻ \oo{Chính sách tài chính}: hạn mức
công nợ và rào giá. Rào giá quyết định báo giá nào phải trình duyệt (mục \ref{dac-thu}). Cách đặt hai số này ở
\kb{kinh-doanh-tai-chinh}{Chính sách tài chính của khách}.

% ============================================================
\Needspace*{16\baselineskip}
\section{Lập phiếu tính giá}\label{tinh-gia}

\begin{boicanh}
Phiếu tính giá ra giá vốn của từng sản phẩm từ quy cách, giấy, máy in và chuỗi công đoạn. Đây là việc của Trợ lý
kinh doanh. Một phiếu có thể nhiều sản phẩm, và một phiếu chỉ ra được một báo giá.

Thứ tự làm: chọn khách, thêm sản phẩm, điền sáu khối của form, bấm \oo{Xong}, bấm \oo{Tính giá \& lưu}, đọc giá
vốn, rồi bấm \oo{Báo giá →}.
\end{boicanh}

Phiếu tính giá không có ô đơn giá nào để gõ tay. Mọi con số đều lấy từ danh mục: giá giấy, máy in, công thức tính
tiền của công đoạn và vật tư. Danh mục thiếu hoặc sai thì giá sai theo. Mỗi mục dưới đây ghi rõ con số lấy từ mục
nào của tài liệu \tl{Hướng dẫn cấu hình hệ thống}, bấm vào tên mục để mở đúng chỗ (hai file PDF phải nằm đúng thư
mục như khi tải về).

\subsection{Mở phiếu mới}\label{tinh-gia-moi}

\rong[0.8\linewidth]{kd-10-ds-phieu.png}

\duong{Kinh doanh > Tính giá}. \so{1} \oo{Lập phiếu tính giá} mở phiếu mới. \so{2} Bấm mã phiếu để mở phiếu cũ.
Cột \oo{Báo giá} cho biết phiếu đã ra báo giá chưa.

\rong[0.8\linewidth]{kd-11-phieu-moi.png}

\so{1} Chọn khách. Chưa chọn khách thì phiếu vẫn lưu và tính được, nhưng nút báo giá bị khoá. Giao tới và người
nhận lấy theo khách, bấm vào để đổi. Khách chưa có thì tạo trước, xem \kb{kinh-doanh-khach}{Khai khách hàng}.

\so{2} \oo{Thêm sản phẩm}: mở form cấu hình (các mục \ref{tinh-gia-kho} đến \ref{tinh-gia-gv}). Hàng khách đã
từng đặt thì gõ tên rồi chọn trong gợi ý. Phần mềm hỏi \oo{Nạp cấu hình tái bản?}, đồng ý thì cả quy cách, giấy,
máy, công đoạn của lần trước được điền sẵn, chỉ cần sửa số lượng.

\so{3} \oo{Tính giá \& lưu}: tính lại và lưu phiếu.

\so{4} \oo{Báo giá →}: lập báo giá từ phiếu này (chương \ref{bao-gia}).

Form sản phẩm có sáu khối bên trái, đánh số 1 đến 6. Bên phải là sơ đồ bình bài, số tờ và giá vốn của sản phẩm,
tự cập nhật ngay khi bạn gõ, chưa cần bấm tính giá.

\Needspace*{20\baselineskip}
\subsection{Khối 1: sản phẩm và khổ chi tiết}\label{tinh-gia-kho}

\canh{0.6}{kd-12-sp-1.png}{%
\so{1} Tên sản phẩm. Tên này in ra báo giá, nên đặt theo cách khách gọi.

\so{2} Số lượng và đơn vị tính (cái, hộp, cuốn, tờ). Đơn vị tính là đơn vị giá vốn chia ra ở cuối: giá vốn
\q{đ/cuốn}, \q{đ/hộp}.

\so{3} \so{4} Dài và rộng chi tiết phẳng, tính bằng mm: khổ trải phẳng của một con trước khi bế, gấp, dán. Hộp
thì là khổ trải hộp, sách thì là khổ một trang.

\so{5} \oo{Số bài in}: số mặt in khác nhau phải ra kẽm riêng. Tờ rơi, hộp, name card là 1. Sách nhiều tay thì
mỗi tay là một bài, bấm \oo{tính từ số trang} để phần mềm đếm.}

Hai ô ngay dưới \oo{Số bài in}:

\begin{itemize}
\item \oo{Bleed (tràn lề)}: phần in tràn ra mỗi cạnh của con để xén không hở viền trắng, thường 3 mm. Gõ 0 nếu
hàng không in tràn lề.
\item \oo{Khe cắt giữa con}: khoảng hở giữa hai con cạnh nhau trên tờ in. Gõ 0 nếu bình sát, cắt chung một nhát.
\end{itemize}

Hai ô này ảnh hưởng thẳng tới số con xếp được trên một tờ in, nên ảnh hưởng tới số tờ và tiền giấy.

\Needspace*{14\baselineskip}
\canh{0.55}{kd-18-so-trang.png}{%
Sách, catalogue, tạp chí in nhiều tay:

\so{1} Bấm \oo{tính từ số trang}.

\so{2} \oo{Số trang}: tổng số trang của phần đang tính, ví dụ ruột 64 trang.

\so{3} \oo{Trang mỗi tay}: số trang một tay sách, hoặc \so{4} bấm nhanh 8, 16, 32 trang.

\so{5} Kết quả: 64 ÷ 16 = 4 tay đủ, tức 4 bài in.

\so{6} \oo{Dùng 4} để điền vào ô \oo{Số bài in}.}

Có số bài in khác 1 thì sơ đồ bên phải đổi thành \oo{Sơ đồ tay sách}, và khối 3 tự chuyển sang kiểu in sách (mục
\ref{tinh-gia-in}). Số kẽm nhân theo số bài in.

Bìa và ruột in khác giấy, khác màu thì làm hai sản phẩm riêng trong cùng phiếu, rồi gộp lại thành một dòng khi báo
giá (mục \ref{tinh-gia-gop}).

\Needspace*{12\baselineskip}
\subsection{Khối 2: giấy nguyên}\label{tinh-gia-giay}

Ô \so{6} ở ảnh trên. Chọn \oo{Giấy} và, nếu là hàng sóng, chọn thêm \oo{Sóng}. Một sản phẩm có tối đa một Giấy và
một Sóng. Gõ \oo{Dài nguyên} và \oo{Rộng nguyên} là khổ tờ giấy mua về, ví dụ 1090 × 790 mm.

Đơn giá hiện ngay dưới ô chọn giấy, ví dụ \q{18.515 đ / kg}. Số này lấy từ danh mục, không sửa trên phiếu. Giá giấy
đổi thì báo người khai danh mục sửa ở \kb{danh-muc-giay}{Giấy}. Ở đó cũng có định lượng và công thức tính tiền giấy
mà phiếu dùng.

Giấy không có trong danh sách nghĩa là chưa khai trong danh mục Giấy.

\Needspace*{20\baselineskip}
\subsection{Khối 3: kỹ thuật in và màu in}\label{tinh-gia-in}

\canh{0.6}{kd-13-sp-2.png}{%
\so{1} \oo{Máy in}. Chọn máy thì phần mềm điền \oo{Khổ tờ in dài}, \oo{Khổ tờ in rộng} và chừa sẵn phần nhíp
giấy, đuôi, lề hông theo thông số máy. Khổ tờ in sửa được, ví dụ khi cắt một tờ nguyên ra hai tờ in. Máy khai ở
\kb{danh-muc-may}{Thiết bị \& Máy móc}.

\so{2} \oo{Quy cách in} (giải thích dưới).

\so{3} \oo{Mực in}: bấm các màu C, M, Y, K, hoặc \oo{+ Pantone} cho màu pha.}

Với hàng không phải sách, ô \oo{Số con} có nút \oo{tự bình bài}: phần mềm thử các cách xếp con lên tờ in (xoay
hoặc không xoay) và lấy cách được nhiều con nhất. Sơ đồ bên phải vẽ cách xếp đó, kèm số con mỗi tờ và hiệu suất
(phần trăm diện tích tờ in dùng được). Gõ tay số con nhỏ hơn nếu xưởng muốn xếp khác.

\oo{Quy cách in} quyết định số kẽm:

\begin{center}
\small
\begin{tabularx}{0.95\linewidth}{@{}l X X@{}}
\toprule
Quy cách & Nghĩa & Số kẽm một bài \\
\midrule
1 mặt & Chỉ in mặt A & Số màu mặt A \\
AB & Hai mặt, mỗi mặt một bộ kẽm & Màu mặt A cộng màu mặt B \\
Tự trở & Hai mặt chung một bộ kẽm, lật tờ theo chiều ngang & Số màu khác nhau của cả hai mặt \\
Trở nhíp & Hai mặt chung một bộ kẽm, lật tờ theo chiều nhíp & Số màu khác nhau của cả hai mặt \\
\bottomrule
\end{tabularx}
\end{center}

Ví dụ mặt A in CMYK, mặt B chỉ in K. In AB ra 4 + 1 = 5 kẽm. In tự trở ra 4 kẽm, vì kẽm K dùng chung cho hai mặt.
Tổng kẽm của sản phẩm bằng số kẽm một bài nhân số bài in.

\Needspace*{14\baselineskip}
\canh{0.55}{kd-17-sach-in.png}{%
Kiểu in sách (số bài in khác 1), ví dụ ruột catalogue 64 trang, 4 tay:

\so{1} Quy cách hiện nhãn \oo{sách in 2 mặt}, không còn \oo{1 mặt}.

\so{2} \oo{Trang mỗi mặt}: nhãn \oo{theo tay} nghĩa là tự lấy bằng trang mỗi tay chia 2 mặt, ở đây 16 ÷ 2 = 8.

\so{3} Màu của mặt A và mặt B.

\so{4} Tổng kẽm 32, ghi rõ ở \so{5}: 4 + 4 = 8 kẽm mỗi tay, nhân 4 tay.

\so{6} Ô của vật tư mực: \oo{Độ phủ mực (\%)}, mục \ref{tinh-gia-cd}.

\so{7} Vật tư \oo{Kẽm CTP} tự lấy đúng 32 tấm.}

\Needspace*{16\baselineskip}
\subsection{Khối 4: chuỗi công đoạn và vật tư}\label{tinh-gia-cd}

Chuỗi công đoạn là các bước làm ra sản phẩm, bắt đầu từ In, theo đúng thứ tự xưởng làm: In, cán màng, bế, dán,
đóng gói. Các bước trước in như cắt tờ, ra kẽm không đưa vào chuỗi. Ở ảnh khối 3:

\begin{itemize}
\item \so{4} Danh sách bước. Bấm một bước để xem vật tư của bước đó ở cột bên phải. Mũi tên lên, xuống đổi thứ tự;
nút thùng rác bỏ bước. \oo{Thêm công đoạn} ở cuối danh sách thêm một bước, gõ tên để tìm.
\item \so{5} \oo{+ Thêm vật tư}: thêm vật tư cho bước đang chọn, gõ tên để tìm. Vật tư mà danh mục công đoạn đã gắn
sẵn thì tự có, không cần thêm.
\end{itemize}

Tiền của một bước gồm hai phần, đều tính bằng công thức có sẵn trong danh mục:

\begin{itemize}
\item Tiền công đoạn, theo công thức tính giá của công đoạn. Xem \kb{danh-muc-cong-doan}{Công đoạn}, tab Công thức
tính giá. Công đoạn chưa có công thức thì dòng này ra 0 đ.
\item Tiền từng vật tư, theo công thức tính giá của vật tư. Lượng vật tư (23,15 kg mực, 32 tấm kẽm, 1,92 kg dung
dịch ở ảnh khối 3) cũng tự tính. Xem \kb{danh-muc-vat-tu}{Vật tư khác}.
\end{itemize}

Một số vật tư cần người tính giá cho thêm một con số mà chỉ nhìn mẫu mới biết. Ô đó hiện ngay dưới tên vật tư, ví
dụ \oo{Độ phủ mực (\%)} của mực in, \oo{Diện tích phủ UV (\%)} của keo UV. Đây là chip riêng của vật tư, khai ở danh
mục Vật tư khác. Chip lạ không hiểu nghĩa thì tra \kb{phu-luc-chip}{Bảng chip trong công thức}.

\Needspace*{14\baselineskip}
\canh{0.6}{kd-15-sp-3.png}{%
Bước nào trong danh mục có bật \oo{Bước này cần khuôn} thì hiện khối \oo{Phí khuôn}:

\so{1} \oo{Có sẵn}: xưởng đã có khuôn của hàng này, không tính tiền.

\so{2} \oo{Làm mới}: gõ \oo{Phí làm khuôn}. Phí tính một lần cho cả lô, không nhân theo số lượng.

\so{3} \so{4} Khối 5 và 6, mục \ref{tinh-gia-chi-phi}.

\so{5} \oo{Xong} để đóng form.}

Phiếu tính giá chỉ hỏi khuôn có sẵn hay làm mới, không chọn khuôn cụ thể trong kho. Chọn khuôn là việc của Kế hoạch
sản xuất khi lên lệnh. Kho khuôn khai ở \kb{danh-muc-khuon}{Khuôn}.

\Needspace*{14\baselineskip}
\subsection{Khối 5 và 6: giao hàng, chi phí khác}\label{tinh-gia-chi-phi}

\canh{0.55}{kd-1a-chi-phi-khac.png}{%
\so{1} \oo{Tổng phí chở hàng}: tiền chở cả lô tới khách. Gõ một lần, phần mềm chia đều theo số lượng.

\so{2} \oo{Khoản chi thêm}: tên khoản và số tiền cho cả lô, ví dụ gấp, bắt cuốn, vào keo gáy thuê ngoài 5.000.000 đ.

\so{3} Số đã chia theo đơn vị: 5.000.000 ÷ 2.000 cuốn = 2.500 đ/cuốn.

\so{4} \oo{+ Thêm khoản chi} cho khoản tiếp theo.

\so{5} Tổng chi phí khác.}

Chi phí khác cộng vào giá vốn nhưng không in thành dòng riêng trên báo giá gửi khách. Việc xưởng không có công đoạn
trong danh mục (thuê ngoài, mua hàng gửi kèm) thì đưa vào đây.

\Needspace*{18\baselineskip}
\subsection{Số tờ máy tự tính}\label{tinh-gia-so-to}

Ô \so{7} ở ảnh khối 1, đọc từ trên xuống:

\begin{itemize}
\item \oo{Thành phẩm cần}: số lượng ở khối 1, ví dụ 2.000 cuốn.
\item \oo{Tờ in cần (chưa hao)}: số lượng chia số con mỗi tờ in, làm tròn lên. Sách thì là số lượng nhân số tay:
2.000 cuốn × 4 tay = 8.000 tờ in.
\item \oo{+ Bù hao công đoạn}: số tờ chạy thêm để bù hàng hỏng. Bấm dòng này để xem hao của từng bước.
\item \oo{= Tờ in}: tờ in cần cộng bù hao. Đây là số tờ đưa vào máy in.
\item \oo{= Tờ nguyên}: số tờ giấy phải mua, bằng tờ in chia số tờ in cắt được từ một tờ nguyên.
\end{itemize}

Bù hao không gõ trên phiếu. Mỗi công đoạn có bảng bù hao riêng theo bậc số lượng, khai ở \kb{danh-muc-cong-doan}{Công
đoạn}, phần Bù hao. Phần mềm cộng ngược từ bước cuối về bước in: bước bế hao 50 tờ, bước cán hao 30 tờ thì máy in
phải ra thêm 80 tờ. Ví dụ trong ảnh có bù hao 0 tờ vì công đoạn In của dữ liệu mẫu chưa khai bù hao. Thấy bù hao 0
mà thực tế xưởng có hao thì báo người khai danh mục.

\Needspace*{18\baselineskip}
\subsection{Giá vốn sản phẩm}\label{tinh-gia-gv}

\canh{0.45}{kd-19-gv-sp.png}{%
Kéo cột phải xuống cuối:

\so{1} \oo{Tờ nguyên}: 8.000 tờ in, khổ tờ in bằng khổ giấy nguyên nên 1 tờ in một tờ nguyên.

\so{2} Nguyên vật liệu: giấy, kẽm, mực, dung dịch. Di chuột vào dấu hỏi cạnh tên để xem phép tính.

\so{3} \oo{Công đoạn}: tiền từng bước, ở đây chỉ có In 6.080.000 đ.

\so{4} \oo{Chi phí khác}.

\so{5} \oo{Giá vốn sản phẩm} 34.743.844 đ, chia 2.000 cuốn ra 17.372 đ/cuốn.}

Xong thì bấm \oo{Xong}, về phiếu bấm \oo{Tính giá \& lưu}.

\Needspace*{18\baselineskip}
\subsection{Đọc kết quả trên phiếu}\label{tinh-gia-ket-qua}

\rong[0.85\linewidth]{kd-14-phieu-kq.png}

\so{1} \oo{Tính giá}: tính lại sau khi sửa.

\so{2} \oo{Nhân bản}: tạo phiếu mới giống hệt. Phiếu đã có báo giá mà khách muốn báo thêm một phương án thì nhân
bản phiếu rồi sửa trên bản sao, vì mỗi phiếu chỉ ra một báo giá.

\so{3} Phiếu đã có báo giá thì nút đổi thành \oo{Mở BG\ldots{} →}.

\so{4} Bảng sản phẩm: số lượng, giá vốn, đơn giá vốn. Bấm tên để sửa cấu hình. Hai nút cuối dòng là nhân bản và
xoá sản phẩm.

\so{5} \oo{Chi tiết dòng giá vốn}: bốn nhóm \oo{Nguyên vật liệu}, \oo{Công đoạn}, \oo{Giao hàng}, \oo{Chi phí khác}.

Khối tối màu bên phải là \oo{Tổng giá vốn nội bộ}. Đây là giá vốn, chưa có lợi nhuận. Lợi nhuận, chiết khấu, kê
đặt ở báo giá.

\Needspace*{14\baselineskip}
\canh{0.55}{kd-1b-dien-giai.png}{%
Mỗi dòng của bảng chi tiết có \so{1} \oo{Số tiền} cả lô, \so{2} \oo{Đ/TP} là tiền trên một thành phẩm, và
\so{3} \oo{Diễn giải} là phép tính ra số đó.

Dòng giấy của ruột catalogue đọc như sau: 0,148 kg/m\textsuperscript{2} (định lượng 148 g) × 1,09 m × 0,79 m (khổ
nguyên) × 8.000 tờ × 18.515 đ/kg = 18.876.827 đ. Chia 2.000 cuốn ra 9.438,41 đ/cuốn.

\so{4} Dòng kẽm: 32 bản × 62.200 đ ÷ 2.000 = 995,20 đ/cuốn.}

Dùng cột Diễn giải để đối chiếu với cách tính tay và để trả lời khi sale hỏi vì sao giá ra như vậy. Sale không thấy
bảng này, chỉ thấy giá vốn tổng. Ai thấy chi tiết giá vốn do quyền của vai quyết định, xem
\kb{kinh-doanh-quyen}{Quyền của sáu màn}.

\Needspace*{18\baselineskip}
\subsection{Sản phẩm nhiều phần: gộp một dòng khi báo giá}\label{tinh-gia-gop}

Catalogue có ruột, bìa, tờ phiếu kẹp giữa. Mỗi phần khác giấy, khác màu, khác công đoạn nên tính thành một sản
phẩm riêng. Nhưng khách chỉ cần thấy một dòng \q{Catalogue A4 64 trang}. Phiếu PTG-2026-0007 dưới đây làm như vậy.

\rong[0.8\linewidth]{kd-16-gop.png}

\so{1} Tick các sản phẩm muốn gộp. Chỉ gộp được các phần cùng số lượng.

\so{2} Thanh \oo{N dòng đã chọn} hiện ra. Gõ tên in ra báo giá, ví dụ Catalogue A4 64 trang.

\so{3} \oo{Gộp khi báo giá}. Phiếu tự lưu.

\so{4} Dòng nhóm: tên gộp, ghi chú \oo{3 sản phẩm gộp 1 dòng khi báo giá}, và ô chọn đơn vị tính in ra báo giá.
Giá vốn nhóm là tổng ba phần: 40.127.653 đ, chia 2.000 ra 20.064 đ/cuốn.

\so{5} \oo{Bỏ gộp} để tách lại.

Gộp chỉ đổi cách in ra báo giá. Các phần vẫn tính riêng, và khi lên lệnh sản xuất vẫn là từng phần.

\Needspace*{18\baselineskip}
\subsection{Số ra lạ: sửa ở đâu}\label{tinh-gia-loi}

Phần lớn số lạ là do danh mục. Trợ lý không tự sửa danh mục mà báo người khai danh mục, kèm mã phiếu và dòng bị lạ.

\begin{center}
\small
\begin{tabularx}{\linewidth}{@{}X X X@{}}
\toprule
Thấy gì trên phiếu & Thường do & Sửa ở \\
\midrule
Dòng giấy 0 đ hoặc lệch tính tay & Đơn giá, định lượng hoặc công thức giấy sai & \kbn{danh-muc-giay}{Giấy} \\
Dòng công đoạn 0 đ & Công đoạn chưa có công thức tính giá & \kbn{danh-muc-cong-doan}{Công đoạn} \\
Vật tư 0 đ, lượng 0 & Vật tư chưa có công thức tính giá hoặc định mức & \kbn{danh-muc-vat-tu}{Vật tư khác} \\
Vật tư hay dùng không tự có ở một bước & Công đoạn chưa gắn vật tư ở tab Vật tư & \kbn{danh-muc-cong-doan}{Công đoạn} \\
Bù hao 0 tờ dù xưởng có hao & Công đoạn chưa khai bù hao & \kbn{danh-muc-cong-doan}{Công đoạn} \\
Không có máy cần chọn, khổ tờ in sai & Máy chưa khai hoặc sai thông số & \kbn{danh-muc-may}{Thiết bị \& Máy móc} \\
Không hiện khối Phí khuôn ở bước bế & Công đoạn chưa bật Bước này cần khuôn & \kbn{danh-muc-cong-doan}{Công đoạn} \\
Không hiểu một chip trong diễn giải & & \kbn{phu-luc-chip}{Bảng chip trong công thức} \\
\bottomrule
\end{tabularx}
\end{center}

Danh mục mới khai xong thì người khai nên kiểm lại bằng một phiếu mẫu, theo \kb{kinh-doanh-kiem}{Kiểm bằng một phiếu
tính giá mẫu}.

% ============================================================
\Needspace*{16\baselineskip}
\section{Báo giá}\label{bao-gia}

\begin{boicanh}
Báo giá lập từ phiếu tính giá, không lập tay. Nút \oo{Báo giá mới} ở màn \duong{Kinh doanh > Báo giá in ấn} cũng
đưa về màn Tính giá để chọn phiếu. Báo giá đi qua các trạng thái: Nháp, (Chờ duyệt, Đã duyệt nếu là báo giá đặc
thù), Đã gửi khách, Khách đồng ý hoặc Từ chối, Đã lên đơn.
\end{boicanh}

\subsection{Đặt giá bán}\label{bao-gia-gia}

Trên phiếu tính giá bấm \oo{Báo giá →}. Báo giá mới mở ra ở trạng thái \oo{Nháp}, mỗi sản phẩm của phiếu là một
dòng.

\rong[0.85\linewidth]{kd-20-bg-nhap.png}

\so{1} Phiếu tính giá nguồn. Giá vốn trên báo giá khoá theo phiếu, muốn đổi giá vốn thì sửa phiếu.

\so{2} \oo{Chiết khấu}: phần chi lại cho khách, gõ theo \oo{\%} hoặc \oo{Số}.

\so{3} \oo{Lợi nhuận}: mặc định 20\% trên giá vốn, gõ theo phần trăm hoặc số tiền. Số tiền lợi nhuận hiện ngay
dưới ô.

\so{4} \oo{Giá bán}: tự ra từ giá vốn, chiết khấu, lợi nhuận. Gõ thẳng giá bán cũng được, lợi nhuận tự tính lại.

\so{5} \oo{Kê}: phần chi lại cho người bên khách, cộng thêm vào giá bán. Cột cuối là \oo{Giá bán sau kê}, là giá
in trên báo giá gửi khách.

\so{6} \oo{Xem bản in}: xem trước bản báo giá gửi khách, in hoặc lưu PDF.

\so{7} \oo{Gửi khách}. Báo giá đặc thù thì nút này là \oo{Trình duyệt} (mục \ref{dac-thu}).

Dưới mỗi sản phẩm có \oo{Sửa diễn giải} (dòng mô tả in trên báo giá) và \oo{Thêm ảnh minh họa}.

\rong[0.85\linewidth]{kd-21-bg-duoi.png}

\so{1} \oo{Điều khoản áp dụng}: in trên bản báo giá, sửa được khi báo giá còn nháp. Bấm \oo{Lưu nháp điều khoản}
sau khi sửa.

\so{2} \oo{Hạn hiệu lực báo giá}: mặc định 30 ngày. Quá hạn mà khách chưa chốt thì báo giá tính là \oo{Hết hiệu lực}.

\so{3} \oo{Giá bán đề xuất}: tổng của báo giá, lợi nhuận bình quân, giá vốn, chiết khấu, lợi nhuận, giá bán chưa
VAT.

\so{4} \oo{VAT}: chọn 0\%, 5\%, 8\% hoặc 10\%.

Phiếu tính giá sửa sau khi đã có báo giá thì báo giá hiện băng vàng. Báo giá còn nháp thì bấm \oo{Cập nhật theo
phiếu}. Báo giá đã gửi khách thì bấm \oo{Tạo phiên bản mới theo phiếu}: bản khách đang cầm vẫn giữ trong
\oo{Lịch sử phiên bản}.

\Needspace*{12\baselineskip}
\subsection{Báo giá đặc thù phải trình duyệt}\label{dac-thu}

Báo giá là đặc thù khi rơi vào một trong ba trường hợp: từ 1 tỷ trở lên (trước VAT), bán dưới giá vốn, hoặc chiết
khấu, lợi nhuận, kê nằm ngoài rào giá của khách. Khi đó báo giá hiện khối \oo{Báo giá đặc thù, cần duyệt} và nút
\oo{Gửi khách} đổi thành \oo{Trình duyệt}.

\begin{enumerate}
\item Bấm \oo{Trình duyệt}. Báo giá sang \oo{Chờ duyệt} và người duyệt nhận thông báo.
\item Người có quyền duyệt (thường là giám đốc) mở báo giá, bấm \oo{Duyệt} hoặc \oo{Từ chối} kèm lý do.
\item Được duyệt thì báo giá sang \oo{Đã duyệt}, bạn bấm \oo{Gửi khách}. Bị từ chối thì bấm \oo{Tạo phiên bản
mới}, sửa giá rồi trình lại.
\end{enumerate}

\Needspace*{14\baselineskip}
\subsection{Gửi khách và ghi kết quả}\label{bao-gia-ket-qua}

Bấm \oo{Gửi khách} sau khi đã gửi bản PDF cho khách qua email hoặc Zalo. Phần mềm không tự gửi email, nút này ghi
lại là đã gửi và khoá bản báo giá.

\rong[0.85\linewidth]{kd-22-bg-gui.png}

\so{1} Trạng thái \oo{Đã gửi khách}. \so{2} Khách không lấy thì bấm \oo{Khách từ chối} và ghi lý do. Muốn báo lại
giá mới thì bấm \oo{Tạo phiên bản mới}. \so{3} Khách đồng ý thì bấm \oo{Khách chốt}.

\canh{0.5}{kd-23-khach-chot.png}{%
Hộp \oo{Khách chốt, chọn sản phẩm} hiện ra.

\so{1} \oo{Chọn tất cả} khi khách lấy cả báo giá.

\so{2} Khách chỉ lấy một phần thì bỏ tích sản phẩm khách không lấy. Tổng tiền dưới cùng tính lại theo phần được
chọn.

\so{3} \oo{Xác nhận khách chốt}.}

\rong[0.85\linewidth]{kd-24-bg-dong-y.png}

\so{1} Báo giá sang \oo{Khách đồng ý}. \so{2} Bấm \oo{Tạo đơn hàng} để lên đơn (chương \ref{don}). Mỗi báo giá chỉ
ra một đơn. Đơn tạo xong thì nút đổi thành \oo{Xem đơn hàng}, trạng thái báo giá là \oo{Đã lên đơn}.

% ============================================================
\Needspace*{16\baselineskip}
\section{Đơn hàng bán}\label{don}

\begin{boicanh}
Đơn tạo từ báo giá mở ở trạng thái \oo{Nháp}. Bạn điền số PO, cọc, ngày giao rồi chốt đơn. Từ lúc chốt, đơn đi
qua năm chặng hiện trên đầu ngăn đơn: \oo{Chốt}, \oo{Cọc}, \oo{Sản xuất}, \oo{Giao}, \oo{Hóa đơn}. Bấm vào một chặng
để xem chi tiết chặng đó.
\end{boicanh}

\subsection{Điền thông tin và chốt đơn}\label{don-chot}

\rong[0.85\linewidth]{kd-30-dh-nhap.png}

\so{1} \oo{Sửa}: mở form thông tin đơn. \so{2} Thanh năm chặng, đơn nháp đứng ở \oo{Chốt}. \so{3} Những thứ còn
thiếu để chốt, mỗi dòng một chữ đỏ. \so{4} \oo{Chốt đơn} chỉ bấm được khi không còn dòng đỏ.

\canh{0.45}{kd-31-dh-sua.png}{%
\so{1} \oo{Số PO khách}: bắt buộc để chốt.

\so{2} \oo{Cọc phải thu}: theo \oo{\%} giá trị đơn hoặc \oo{Số} tiền. Để trống là đơn không cần cọc.

\so{3} \oo{Ngày giao cam kết}: bắt buộc để chốt. Gõ ngày hoặc chọn trên lịch.

\so{4} \oo{Lưu ý giao hàng}: dặn người giao, ví dụ giờ giao, gọi trước.

\so{5} \oo{Lưu ý sản xuất}: dặn xưởng, ví dụ canh màu theo mẫu, cấn bế kỹ.

\so{6} \oo{Hàng gấp (ưu tiên sản xuất)}: đánh dấu đơn gấp cho xưởng.

\so{7} \oo{Lưu}.

Địa chỉ giao và người nhận lấy từ báo giá. Muốn đổi thì sửa ở báo giá hoặc hồ sơ khách trước khi lên đơn.}

Lưu xong, ô \oo{Chốt đơn} báo \oo{Đủ điều kiện chốt}. Bấm \oo{Chốt đơn}. Đơn sang \oo{Đã chốt}. Đơn đã chốt không
sửa được nữa, chỉ huỷ được bằng nút \oo{Hủy đơn} kèm lý do.

\Needspace*{14\baselineskip}
\subsection{Chờ cọc}\label{don-coc}

\rong[0.85\linewidth]{kd-32-dh-coc.png}

\so{1} Đơn có cọc thì sau khi chốt đứng ở chặng \oo{Cọc}. \so{2} Dòng việc tiếp theo nói rõ đơn đang chờ ai, còn
thiếu bao nhiêu. \so{3} Thiếu cọc thì đơn chưa xuống sản xuất.

Thu cọc là việc của Kế toán: nút \oo{Lập phiếu thu cọc} chỉ hiện với người có quyền thu tiền. Sale báo khách chuyển
khoản rồi báo Kế toán. Kế toán thu đủ cọc thì đơn tự xuống hàng chờ của Kế hoạch sản xuất, không ai phải bấm
chuyển. Đơn không cần cọc thì xuống hàng chờ ngay khi chốt.

\Needspace*{14\baselineskip}
\subsection{Đơn đang sản xuất}\label{don-sx}

\rong[0.85\linewidth]{kd-33-dh-sx.png}

Chặng \oo{Sản xuất} liệt kê từng sản phẩm của đơn. \so{1} Mã lệnh sản xuất, bấm để mở hồ sơ lệnh. Cột \oo{Đang ở}
cho biết lệnh đang ở bước nào hoặc ở nhà gia công ngoài, cột \oo{Đã có hàng} là số thành phẩm kho đã nhận.

\so{2} Dòng tóm tắt, ví dụ \q{2 lệnh chưa xuống xưởng}: Kế hoạch đã có lệnh nhưng chưa phát hành. \so{3} \oo{Mở
Kế hoạch SX} để xem chi tiết từng công đoạn (cần quyền xem Kế hoạch sản xuất).

Khi lệnh đã chạy xong các công đoạn, dòng tóm tắt nói hàng đang ở khâu KCS hay kho:

\begin{itemize}
\item \q{Chờ KCS gửi thành phẩm sang kho}: hàng đã làm xong, đang chờ KCS kiểm rồi gửi sang kho. Báo khách hàng
đang kiểm lần cuối. Dòng này đứng lâu thì hỏi người làm KCS.
\item \q{Kho chưa nhận \ldots{} KCS đã gửi}, ví dụ \q{Kho chưa nhận 1.000 hộp Hộp thuốc KCS đã gửi}: KCS kiểm xong
và đã gửi, thủ kho chưa nhận vào kho. Hỏi thủ kho.
\end{itemize}

Hạn giao cam kết ở bên phải dòng tóm tắt, quá hạn thì có chữ đỏ \q{quá hạn giao}.

% ============================================================
\Needspace*{16\baselineskip}
\section{Trả lời khách: hàng tới đâu rồi}\label{theo-doi}

\begin{boicanh}
\duong{Sản xuất > Theo dõi sản xuất} gom mọi đơn đã xuống xưởng vào một bảng. Đây là màn trả lời khách nhanh nhất:
mỗi dòng một đơn, thấy ngay đơn đang làm gì, có vướng gì, còn bao lâu tới hạn. Bảng chỉ hiện đơn trong phạm vi quyền của bạn.
\end{boicanh}

Các tab trên cùng lọc nhanh: \oo{Cần để ý} (trễ hạn, đang vướng), \oo{Đang làm}, \oo{Chờ giao}, \oo{Đã giao đủ}.
Cột \oo{Đang vướng} ghi lý do, ví dụ \q{Đã quá hạn giao}, \q{Thiếu vật tư}. Bấm một dòng để mở chi tiết.

\rong[0.85\linewidth]{kd-40-theo-doi.png}

\so{1} Từng mặt hàng: số đặt, số đã nhập kho, số đang có hàng, số đã giao.

\so{2} Lệnh sản xuất của mặt hàng và bước đang làm. Ví dụ trên là lệnh gia công trọn gói ở nhà gia công ngoài.

\so{3} \oo{Đang vướng}: lý do đơn chậm. \q{KCS không đạt} nghĩa là có lần KCS kết luận cả lượt kiểm không đạt.
Hàng lỗi bị trừ khỏi số đạt nên đơn có thể thiếu, Kế hoạch sản xuất quyết định làm bù hay không. Hỏi Kế hoạch sản xuất trước khi hẹn lại ngày giao
với khách.

\so{4} \oo{Giao hàng}: còn bao nhiêu chưa xuất hàng.

Muốn xem sâu hơn từng bước của một lệnh thì mở \duong{Sản xuất > Hồ sơ lệnh sản xuất}, tìm theo mã lệnh hoặc mã đơn.

% ============================================================
\Needspace*{16\baselineskip}
\section{Xuất hàng}\label{giao}

\begin{boicanh}
Kinh doanh bấm \oo{Xuất hàng} ở chặng \oo{Giao} của đơn. Chỉ xuất được phần kho đã nhận. Phần mềm gửi đề nghị xuất
sang kho, thủ kho lập phiếu xuất và ghi sổ. Xuất thiếu thì bấm thêm lần sau cho phần còn lại.
\end{boicanh}

\rong[0.85\linewidth]{kd-50-giao.png}

\so{1} Bấm chặng \oo{Giao} trên thanh chặng của đơn.

\so{2} Bảng từng mặt hàng: khách đặt, đã giao, chờ kho xuất, còn phải giao, \oo{Đang có ở kho} (số đang nằm ở
Kho thành phẩm, kể cả phần đã giữ cho lần xuất đang chờ) và \oo{Xuất được ngay}. Xuất được ngay là số nhỏ hơn giữa
còn phải giao và phần ở kho chưa ai giữ, nên có thể bằng 0 dù kho vẫn có hàng.

\so{3} Dòng nhắc: xuất được tối đa bao nhiêu. Chưa xuất được gì thì nút bị mờ và dòng nhắc nói rõ từng trường hợp:
kho chưa có mặt hàng nào và đang chờ lệnh nào, hoặc mặt hàng nào có ở kho nhưng đã giữ hết cho lần xuất đang chờ
hay cho đơn khác của khách.

\so{4} \oo{Xuất hàng}: mở form gồm \oo{Hàng xuất đợt này} (số lượng từng mặt hàng, không quá số giao được),
\oo{Ngày cần giao}, \oo{Giao tới}, \oo{Người nhận}, \oo{Ghi chú cho kho}. Bấm \oo{Gửi kho xuất}.

\so{5} Sau khi gửi, lần xuất hiện trong bảng lần xuất của chặng Giao, mỗi mặt hàng một dòng, cột \oo{Kho} ghi
\q{Chờ kho xuất}. Thủ kho thấy đề nghị ở \duong{Kho hàng > Yêu cầu nhập xuất}, lập phiếu xuất rồi ghi sổ. Cột
\oo{Lần xuất} ghi số phiếu xuất của kho (PXK…), chưa có phiếu thì ghi \q{chờ kho lập phiếu}. Ghi sổ xong, cột
\oo{Kho} đổi thành \q{Đã xuất} kèm giờ xuất, nút \oo{In phiếu xuất} mở ra, và số kho xuất được tính là đã giao
cho khách. Bảng hoá đơn gọi lần xuất bằng đúng số phiếu này.

% ============================================================
\Needspace*{16\baselineskip}
\section{Hoá đơn, tiền về và hoa hồng}\label{tien}

\begin{boicanh}
Kế toán ghi hoá đơn theo phần đã giao và thu tiền. Kinh doanh không làm hai việc này nhưng cần đọc được: đơn đã ghi
hoá đơn chưa, khách còn nợ bao nhiêu, hoa hồng của mình ra sao.
\end{boicanh}

\rong[0.85\linewidth]{kd-52-hoa-don.png}

\so{1} Chặng \oo{Hóa đơn} của đơn. \so{2} \oo{Tổng đơn}, \oo{Đã giao}, \oo{Đã giao chưa ghi} (hàng kho đã xuất
mà Kế toán chưa ghi hoá đơn), \oo{Đã ghi hóa đơn}, \oo{Còn phải thu}. Có hoá đơn thì khách mới phát sinh công nợ.

Cột \oo{Cọc} và \oo{Giao khách} ở danh sách \duong{Kinh doanh > Đơn hàng bán} cho thấy nhanh đơn nào đủ cọc, giao
tới đâu. Khách nợ quá hạn hay vượt hạn mức thì Kế toán thấy ở màn Công nợ phải thu.

\subsection{Hoa hồng}\label{hoa-hong}

Hoa hồng thuộc người phụ trách khách, ghi trên đơn ở dòng \q{\ldots{} phụ trách}. Phần mềm tính như sau:

\begin{itemize}
\item Hoa hồng chỉ phát sinh khi tiền về, tức khi Kế toán lập phiếu thu. Đơn chốt, giao, ghi hoá đơn mà khách chưa
trả thì chưa có hoa hồng.
\item Tiền tính hoa hồng là tiền thu về trước VAT, trừ chiết khấu và kê của đơn.
\item Phần trăm tra theo biểu bậc của bạn, theo số ngày từ ngày hoá đơn tới ngày thu tiền: khách trả càng sớm phần
trăm càng cao. Tiền cọc thu trước khi có hoá đơn tính theo bậc đầu tiên.
\item Biểu bậc chụp vào đơn lúc chốt. Biểu đổi về sau thì đơn đã chốt vẫn tính theo biểu cũ.
\end{itemize}

Biểu bậc do bộ phận nhân sự khai trong mức lương của từng người. Hoa hồng của tháng hiện trên phiếu lương, dòng
\oo{Hoa hồng kinh doanh}: \duong{Nhân sự \& Lương > Lương}, tab \oo{Phiếu lương của tôi}. Thắc mắc số nào thì hỏi
bộ phận lương: họ mở được bảng đối chiếu từng lần tiền về của bạn, mỗi dòng ghi số tiền, số ngày và phần trăm đã áp.

% ============================================================
\Needspace*{16\baselineskip}
\section{Câu hỏi hay gặp}\label{hoi-dap}

\begin{description}[style=nextline,leftmargin=0pt,font=\normalfont\tieude\small]
\item[Nút Báo giá trên phiếu tính giá bị mờ.]
Phiếu chưa chọn khách hoặc chưa lưu. Chọn khách, bấm \oo{Tính giá \& lưu} rồi bấm lại.

\item[Khách muốn báo hai phương án cho cùng một hàng.]
Mỗi phiếu ra một báo giá. Bấm \oo{Nhân bản} trên phiếu, sửa phương án hai trên bản sao, rồi ra báo giá thứ hai.

\item[Khách đòi giảm giá sau khi đã gửi báo giá.]
Báo giá đã gửi bị khoá. Bấm \oo{Khách từ chối} với lý do \q{xin giảm giá}, rồi \oo{Tạo phiên bản mới}, sửa giá và gửi
lại. Các bản cũ xem ở \oo{Lịch sử phiên bản}, nút \oo{So sánh} đặt hai bản cạnh nhau.

\item[Không thấy nút Gửi khách, chỉ có Trình duyệt.]
Báo giá đặc thù (mục \ref{dac-thu}). Trình duyệt, chờ giám đốc duyệt rồi mới gửi khách được.

\item[Không chốt được đơn.]
Đọc các dòng đỏ trong ô \oo{Chốt đơn}: thường thiếu số PO hoặc ngày giao cam kết. Bấm \oo{Sửa} để điền.

\item[Đơn đã chốt mà xưởng chưa thấy.]
Đơn có cọc mà Kế toán chưa thu đủ. Chặng \oo{Cọc} ghi còn thiếu bao nhiêu.

\item[Nút Xuất hàng bị mờ.]
Kho chưa nhận thành phẩm của đơn, hoặc hàng ở kho đã giữ hết cho lần xuất đang chờ. Dòng nhắc dưới bảng nói
rõ trường hợp nào, đang chờ lệnh nào. Cột \oo{Đang có ở kho} cho biết kho thực còn bao nhiêu.

\item[Khách của mình đứng tên người khác.]
Báo trưởng phòng dùng chức năng \oo{Điều chuyển} ở màn Khách hàng để chuyển khách về đúng người, xem
\kb{kinh-doanh-dieu-chuyen}{Đổi người phụ trách}.
\end{description}

\end{document}
